% theory_market_impact.tex
% BTC price impact, reflexivity amplification, and systemic risk

\section{BTC Price Impact and Reflexivity}
\label{sec:market_impact}

Strategy's BTC acquisitions are large relative to daily trading volume.
This section models the price impact of purchases, derives a reflexivity
amplification factor, and analyzes systemic risk from corporate BTC
treasury adoption.

\subsection{Price Impact Model}

We augment the BTC price dynamics (Section~2.2) with a permanent
price impact term:
\begin{equation}
  \frac{dS_t}{S_t} = \mu_S\,dt + \sigma_S\,dW_t^{S,\mathbb{P}}
  + \eta\,\frac{dH_t}{V_t},
  \label{eq:price_impact}
\end{equation}
where:
\begin{itemize}
  \item $V_t$ is the daily BTC trading volume (in BTC),
  \item $\eta > 0$ is the Kyle lambda---the permanent price impact
    per unit of order flow normalized by volume,
  \item $dH_t / V_t$ is Strategy's purchase flow as a fraction
    of market volume.
\end{itemize}

\begin{remark}
This follows the Kyle (1985) linear price impact framework.
The normalization by $V_t$ makes $\eta$ dimensionless and
captures the intuition that the same dollar purchase has less
impact in a liquid market.
For BTC, empirical estimates suggest $\eta \in [0.1, 0.5]$
depending on the execution horizon and venue.
\end{remark}

Substituting the holding dynamics $dH_t = \alpha\pi_t^+\,dt + dH_t^{\mathrm{jump}}$:
\begin{equation}
  \frac{dS_t}{S_t} = \left(\mu_S + \eta\,\frac{\alpha\pi_t^+}{V_t}\right)dt
  + \sigma_S\,dW_t^{S,\mathbb{P}}
  + \eta\,\frac{dH_t^{\mathrm{jump}}}{V_t}.
  \label{eq:price_impact_expanded}
\end{equation}

The drift of BTC price is now \emph{endogenously increased} when
the premium is positive, creating a feedback channel.

\subsection{Reflexivity and the Amplification Factor}

The feedback loop operates as follows:
\begin{enumerate}
  \item Strategy purchases BTC ($dH_t > 0$).
  \item BTC price rises due to price impact ($dS_t > 0$).
  \item NAV increases: $\mathrm{NAV}_t = H_t S_t - D_t$.
  \item If $\pi_t$ is sticky or increases, equity value $E_t = e^{\pi_t}\mathrm{NAV}_t$ rises.
  \item Higher equity value enables more ATM issuance or convertible placement.
  \item Proceeds fund additional BTC purchases. Return to step 1.
\end{enumerate}

\begin{definition}[Feedback Gain (Market Impact)]
\label{def:feedback_gain_mi}
The \emph{instantaneous feedback gain} is
\begin{equation}
  G_t = \eta\,\frac{\alpha}{V_t}\,
  \frac{\partial\,\text{Issuance Capacity}}{\partial S_t}\,
  \frac{1}{S_t},
  \label{eq:feedback_gain_mi}
\end{equation}
measuring the fraction of a BTC price increase that is recycled
into additional purchases through the financing flywheel.
\end{definition}

To derive $G_t$ explicitly, note that ATM issuance capacity is
proportional to $E_t = e^{\pi_t}(H_t S_t - D_t)$, so:
\begin{align}
  \frac{\partial E_t}{\partial S_t}
  &= e^{\pi_t} H_t + e^{\pi_t}(H_t S_t - D_t)\frac{\partial\pi_t}{\partial S_t}
  \notag\\
  &= E_t\left(\frac{H_t}{H_t S_t - D_t}
    + \frac{\partial\pi_t}{\partial\log S_t}\right)
  = E_t\left(\mathrm{ILE}_t + \gamma_{\pi S}\right)\frac{1}{S_t}
  \notag\\
  &= E_t \cdot \beta_{\mathrm{TOT}}(t) \cdot \frac{1}{S_t}.
  \label{eq:equity_sensitivity}
\end{align}

If a fraction $\phi_t$ of the equity value increase is deployed
to purchase BTC (the ``reinvestment rate''), then:
\begin{equation}
  G_t = \eta\,\frac{\phi_t\,\beta_{\mathrm{TOT}}(t)\,E_t}{V_t\,S_t^2}.
  \label{eq:gain_explicit}
\end{equation}

\begin{proposition}[Reflexivity Amplification Factor]
\label{prop:reflexivity}
The steady-state amplification factor for the BTC price response
to an exogenous demand shock is
\begin{equation}
  R_t = \frac{1}{1 - G_t},
  \label{eq:amplification}
\end{equation}
provided $G_t < 1$ (stability condition).
An exogenous BTC demand shock of $\Delta S^{\mathrm{exog}}$
produces an equilibrium price response of
$\Delta S^{\mathrm{eq}} = R_t \cdot \Delta S^{\mathrm{exog}}$.
\end{proposition}

\begin{proof}
Consider a unit exogenous price increase $\delta$.
Through the feedback loop, this generates additional BTC purchases
of $G_t \delta / \eta$ in volume-normalized terms, producing a
second-round price impact of $G_t \delta$.
Iterating, the total impact is
$\delta(1 + G_t + G_t^2 + \cdots) = \delta/(1 - G_t)$
when $|G_t| < 1$.
\end{proof}

\begin{proposition}[Instability Condition]
\label{prop:instability}
The reflexive feedback loop becomes unstable ($R_t \to \infty$) when
\begin{equation}
  G_t \ge 1
  \quad\Longleftrightarrow\quad
  \eta\,\phi_t\,\beta_{\mathrm{TOT}}(t)\,E_t \ge V_t\,S_t^2.
  \label{eq:instability}
\end{equation}
This is more likely when:
\begin{enumerate}
  \item Price impact $\eta$ is high (illiquid market).
  \item Reinvestment rate $\phi_t$ is high (aggressive issuance).
  \item Total equity elasticity $\beta_{\mathrm{TOT}}(t)$ is high
    (high leverage or pro-cyclical premium).
  \item Equity capitalization $E_t$ is large relative to BTC market
    depth $V_t S_t$.
\end{enumerate}
\end{proposition}

\begin{remark}[Current Calibration]
With $E_t \approx \$16B$ (at $\pi_0 = 0.242$, NAV $\approx \$48B$,
$E_t = e^{0.242}\times\$48B\approx\$61B$---but the negative $\gamma_{\pi S} = -3.60$
yields $\beta_{\mathrm{TOT}} = -2.11$, which is negative),
the feedback gain is in fact \emph{negative} at current calibration.
This reflects the empirical finding that the premium \emph{compresses}
when BTC rises ($\gamma_{\pi S} < 0$), acting as a natural dampener.
The reflexive amplification is thus $R_t < 1$: price impact is
\emph{attenuated} rather than amplified under current market conditions.

Instability would require $\gamma_{\pi S}$ to turn positive
(premium expanding with BTC price), which could occur during
speculative episodes. The November--December 2024 rally, where
$\pi_t$ exceeded 2.0 coincident with BTC's rise to \$100K+,
may represent such an episode.
\end{remark}

\subsection{Discrete Financing Events and Jump Impact}

Large discrete purchases have concentrated price impact.
At event time $\tau_k$, Strategy purchases $\Delta H_{\tau_k}$ BTC,
producing a price jump:
\begin{equation}
  \frac{\Delta S_{\tau_k}}{S_{\tau_k^-}}
  = \eta\,\frac{\Delta H_{\tau_k}}{V_{\tau_k}}.
  \label{eq:jump_impact}
\end{equation}

\begin{theorem}[Expected Price Impact of Financing Program]
\label{thm:expected_impact}
The expected cumulative price impact of Strategy's BTC acquisition
program over horizon $[0,T]$ is
\begin{equation}
  \mathbb{E}^{\mathbb{P}}\!\left[
    \int_0^T \eta\,\frac{dH_u}{V_u}
  \right]
  = \eta\left(
    \frac{\alpha\,\mathbb{E}[\pi_t^+]}{\bar{V}}\,T
    + \lambda_M\,\frac{\mathbb{E}[\Delta H]}{\bar{V}}\,T
  \right),
  \label{eq:cumulative_impact}
\end{equation}
where $\bar{V}$ is the average daily volume.
With calibrated parameters ($\alpha = 0$,
$\lambda_M = 19.05$ events/year,
$\mathbb{E}[\Delta H] = 7{,}044$ BTC,
$\bar{V} \approx 25{,}000$ BTC/day on major exchanges), this gives
an annual price contribution of approximately
$\eta \times 19.05 \times 7{,}044 / 25{,}000 \approx 5.37\eta$,
or $0.54$--$2.68\%$ for $\eta \in [0.1, 0.5]$.
\end{theorem}

\subsection{Systemic Risk from Corporate BTC Adoption}
\label{subsec:systemic}

Strategy holds approximately 761{,}068 BTC out of roughly 19.8M
circulating supply (3.84\%). If other corporations adopt similar
treasury strategies, concentration risk increases.

\begin{definition}[Aggregate Corporate Holdings]
Let $\mathcal{F} = \{1, \ldots, F\}$ be the set of firms holding BTC
as treasury assets. Define aggregate corporate holdings
\begin{equation}
  H_t^{\mathrm{agg}} = \sum_{f \in \mathcal{F}} H_t^{(f)},
  \qquad
  \omega_t = \frac{H_t^{\mathrm{agg}}}{\bar{H}^{\mathrm{supply}}},
  \label{eq:aggregate_holdings}
\end{equation}
where $\bar{H}^{\mathrm{supply}}$ is the effective circulating supply
and $\omega_t$ is the corporate concentration ratio.
\end{definition}

\begin{proposition}[Fire-Sale Externality]
\label{prop:fire_sale}
Consider $F$ firms each holding $H^{(f)}$ BTC with leverage
$\mathrm{ILE}^{(f)}$. If firm $f$ faces a covenant breach or
liquidity crisis when $S_t < S^{(f)}_{\mathrm{distress}}$, it must
liquidate $\Delta H^{(f)}_{\mathrm{sell}} \le H^{(f)}$ BTC.
The aggregate fire-sale impact on BTC price is
\begin{equation}
  \frac{\Delta S^{\mathrm{fire}}}{S_t}
  = -\eta_{\mathrm{sell}}\,
  \frac{\sum_{f: S_t < S^{(f)}_{\mathrm{distress}}} \Delta H^{(f)}_{\mathrm{sell}}}{V_t},
  \label{eq:fire_sale}
\end{equation}
where $\eta_{\mathrm{sell}}$ is the sell-side price impact (potentially
larger than $\eta$ due to adverse selection in distressed sales).
\end{proposition}

This connects to the fire-sale literature of \citet{shleifer1997limits}
and \citet{shleifer2011fire}. The key distinction is that
BTC is a single homogeneous asset with a unified global order book,
so fire sales by any holder affect all holders simultaneously---there
is no asset-class diversification to absorb the shock.

\begin{theorem}[Contagion Cascade Condition]
\label{thm:contagion}
A contagion cascade occurs when the fire sale by one firm pushes
$S_t$ below the distress threshold of another firm. Formally,
order the firms by distress threshold:
$S^{(1)}_{\mathrm{distress}} > S^{(2)}_{\mathrm{distress}} > \cdots$.
A cascade starting from firm $1$ propagates to firm $j$ if
\begin{equation}
  S_t - \eta_{\mathrm{sell}}\,S_t\,
  \frac{\sum_{i=1}^{j-1} \Delta H^{(i)}_{\mathrm{sell}}}{V_t}
  < S^{(j)}_{\mathrm{distress}}.
  \label{eq:contagion}
\end{equation}
The cascade is self-limiting if and only if there exists $j^* < F$
such that the post-liquidation price exceeds all remaining thresholds:
\begin{equation}
  S_t\left(1 - \eta_{\mathrm{sell}}\,
  \frac{\sum_{i=1}^{j^*} \Delta H^{(i)}_{\mathrm{sell}}}{V_t}\right)
  > S^{(j^*+1)}_{\mathrm{distress}}.
  \label{eq:cascade_limit}
\end{equation}
\end{theorem}

\begin{proof}
The proof follows by induction on the ordered firm sequence.
If the post-liquidation price after firm $j-1$ sells is below
$S^{(j)}_{\mathrm{distress}}$, firm $j$ is forced to sell,
further depressing the price by
$\eta_{\mathrm{sell}}\,\Delta H^{(j)}_{\mathrm{sell}} / V_t$.
The cascade terminates at $j^*$ if the cumulative price impact
is insufficient to breach the next threshold.
\end{proof}

\begin{remark}[Current Systemic Assessment]
Strategy is currently the dominant corporate BTC holder.
The nearest comparables (Marathon Digital, Riot Platforms, etc.)
hold 1--2 orders of magnitude less BTC.
Under current concentration, the fire-sale externality is
primarily a single-firm risk rather than a systemic cascade.

However, Strategy's success has catalyzed adoption: at least 70+
public companies now hold BTC on balance sheet.
If aggregate corporate holdings reach 10--15\% of circulating supply
(currently $\sim$5\% including miners and ETFs), the contagion cascade
mechanism becomes material.

Strategy's own distress threshold, with $D_t = \$8.22$B and
$H_t = 761{,}068$ BTC, occurs at
$S_{\mathrm{distress}} = D_t / H_t \approx \$10{,}800$.
Including preferred liquidation claims (\$10.23B),
the full-stack distress threshold is
$(D_t + \text{Pref}_t) / H_t \approx \$24{,}200$.
Both are well below the current price of \$73{,}909, implying
substantial buffer against forced liquidation.
\end{remark}

\subsection{Volume-Adjusted Impact Over Time}

As BTC market depth grows, the price impact of a fixed-size purchase
diminishes. We model volume growth as:
\begin{equation}
  V_t = V_0\,e^{g_V t},
  \label{eq:volume_growth}
\end{equation}
where $g_V$ is the annual growth rate of BTC trading volume.
The time-varying amplification factor becomes:
\begin{equation}
  R_t = \frac{1}{1 - G_0\,e^{(g_E - 2g_S - g_V)t}},
  \label{eq:time_varying_R}
\end{equation}
where $g_E$ is the growth rate of equity value and $g_S$ the growth
rate of BTC price. The feedback loop naturally weakens over time
if BTC volume and price grow faster than Strategy's equity
capitalization---a plausible scenario as the BTC market matures.

\begin{corollary}[Asymptotic Stability]
\label{cor:asymptotic}
If $2g_S + g_V > g_E$, then $G_t \to 0$ and $R_t \to 1$ as
$t \to \infty$: the reflexive feedback loop asymptotically vanishes.
Strategy's price impact becomes negligible as the BTC market
grows relative to any single corporate buyer.
\end{corollary}